% 第 6 题程序框图（本项目首例流程图类图元）：按固定源 PDF 矢量路径换算（1 pt → 0.03515 cm）。
% 原点取图左下角（源页 (154.98, 500.84) pt），y 向上。
% 结构：自上而下主干 + 右侧回环（b←b−2 在右上、a←a+4 在右中）；判断支路标签为拉丁 N/Y（照源，非「否/是」）。
% 节点：圆角矩形 2（起止框，圆角半径 3.99）、矩形 4、菱形 1（半宽 34.82、半高 8.64）、平行四边形 1（斜移 6.13）。
% 8 条连线（含回环折线）+ 8 个实心箭头（2.99×3.98，尖点在目标端；源内线段在箭头前收笔，照源）。
\begin{tikzpicture}[line cap=round,line join=round,font=\small,x=0.03515cm,y=0.03515cm]
% ===== 连线（8 条）=====
\draw (34.81,171.83)--(34.81,161.25);   % 开始 → a←1
\draw (34.81,143.50)--(34.81,132.92);   % a←1 → b←9
\draw (34.81,114.65)--(34.81,84.23);    % b←9 → 菱形（回接点在上）
\draw (69.91,71.73)--(80.49,71.73);     % 菱形 N 支 → a←a+4
\draw (110.52,79.24)--(110.52,89.82);   % a←a+4 → b←b−2（自下而上）
\draw (85.12,101.19)--(38.40,101.19);   % b←b−2 → 主干（自右向左回接）
\draw (34.81,62.80)--(34.81,52.22);     % 菱形 Y 支 → 输出 a
\draw (34.82,31.38)--(34.81,20.80);     % 输出 a → 结束
% ===== 节点（8 个）=====
\draw[rounded corners=3.99pt] (20.82,172.03)--(48.81,172.03)--(48.81,188.86)--(20.82,188.86)--cycle;  % 开始
\draw (17.75,143.70)--(51.88,143.70)--(51.88,157.46)--(17.75,157.46)--cycle;   % a←1
\draw (18.28,114.85)--(51.35,114.85)--(51.35,129.14)--(18.28,129.14)--cycle;   % b←9
\draw (85.32,93.61)--(135.73,93.61)--(135.73,108.77)--(85.32,108.77)--cycle;   % b←b−2
\draw (69.62,71.73)--(34.81,80.37)--(0,71.73)--(34.81,63.09)--cycle;            % 菱形 a>b
\draw (84.27,64.41)--(136.77,64.41)--(136.77,79.04)--(84.27,79.04)--cycle;     % a←a+4
\draw (10.16,31.58)--(16.29,48.43)--(59.46,48.43)--(53.33,31.58)--cycle;       % 平行四边形 输出 a
\draw[rounded corners=3.99pt] (20.82,0)--(48.81,0)--(48.81,17.01)--(20.82,17.01)--cycle;  % 结束
% ===== 实心箭头（8 个，2.99×3.98）=====
\fill (34.81,157.46)--(33.32,161.44)--(36.30,161.44)--cycle;
\fill (34.81,129.14)--(33.32,133.12)--(36.30,133.12)--cycle;
\fill (34.81,80.37)--(33.32,84.35)--(36.30,84.35)--cycle;
\fill (110.52,93.61)--(109.03,89.63)--(112.01,89.63)--cycle;
\fill (34.81,48.43)--(33.32,52.41)--(36.30,52.41)--cycle;
\fill (34.81,17.01)--(33.32,20.99)--(36.30,20.99)--cycle;
\fill (84.27,71.73)--(80.29,70.24)--(80.29,73.22)--cycle;
\fill (34.81,101.19)--(38.79,99.70)--(38.79,102.68)--cycle;
% ===== 标签（照源字框中心，显式坐标）=====
\node[anchor=center,inner sep=0] at (34.75,180.44) {开始};
\node[anchor=center,inner sep=0] at (34.60,150.64) {\(a\leftarrow 1\)};
\node[anchor=center,inner sep=0] at (35.01,121.87) {\(b\leftarrow 9\)};
\node[anchor=center,inner sep=0] at (110.58,101.56) {\(b\leftarrow b-2\)};
\node[anchor=center,inner sep=0] at (34.95,71.71) {\(a>b\)};
\node[anchor=center,inner sep=0] at (69.90,78.99) {N};
\node[anchor=center,inner sep=0] at (110.37,71.90) {\(a\leftarrow a+4\)};
\node[anchor=center,inner sep=0] at (42.44,55.72) {Y};
\node[anchor=center,inner sep=0] at (34.81,40.01) {输出 \(a\)};
\node[anchor=center,inner sep=0] at (34.83,8.50) {结束};
\end{tikzpicture}
